\section{Regular rings and global dimension}
\label{section-regular-finite-gl-dim}

\noindent
We can use the material on rings of finite global dimension
to give another characterization of regular local rings.

\begin{proposition}
\label{proposition-regular-finite-gl-dim}
Let $R$ be a regular local ring of dimension $d$.
Every finite $R$-module $M$ of depth $e$ has a finite free
resolution
$$
0 \to F_{d-e} \to \ldots \to F_0 \to M \to 0.
$$
In particular a regular local ring has global dimension $\leq d$.
\end{proposition}

\begin{proof}
The first part holds in view of Lemma \ref{lemma-regular-mcm-free}
and Lemma \ref{lemma-mcm-resolution}. The last part follows from this
and Lemma \ref{lemma-finite-gl-dim}.
\end{proof}

\begin{lemma}
\label{lemma-finite-gl-dim-primes}
Let $R$ be a Noetherian ring. Let $n \geq 0$ be an integer.
Then $R$ has finite global dimension $\leq n$ if and
only if for all maximal ideals $\mathfrak m$ of $R$
the ring $R_{\mathfrak m}$ has global dimension $\leq n$.
\end{lemma}

\begin{proof}
We saw, Lemma \ref{lemma-localize-finite-gl-dim}
that if $R$ has finite global dimension $n$,
then all the localizations $R_{\mathfrak m}$
have finite global dimension at most $n$.
Conversely, suppose that all the $R_{\mathfrak m}$
have global dimension $\leq n$. Let $M$ be a finite
$R$-module. Let
$0 \to K_n \to F_{n-1} \to \ldots \to F_0 \to M \to 0$
be a resolution with $F_i$ finite free.
Then $K_n$ is a finite $R$-module.
According to
Lemma \ref{lemma-independent-resolution}
and the assumption all the modules $K_n \otimes_R R_{\mathfrak m}$
are projective. Hence by
Lemma \ref{lemma-finite-projective}
the module $K_n$ is finite projective.
\end{proof}

\begin{lemma}
\label{lemma-length-resolution-residue-field}
Suppose that $R$ is a Noetherian local ring
with maximal ideal $\mathfrak m$ and
residue field $\kappa$. In this case
the projective dimension of $\kappa$ is
$\geq \dim_\kappa \mathfrak m / \mathfrak m^2$.
\end{lemma}

\begin{proof}
Let $x_1 , \ldots, x_n$ be elements of $\mathfrak m$
whose images in $\mathfrak m / \mathfrak m^2$ form a basis.
Consider the {\it Koszul complex} on $x_1, \ldots, x_n$.
This is the complex
$$
0 \to \wedge^n R^n \to \wedge^{n-1} R^n \to \wedge^{n-2} R^n \to
\ldots \to \wedge^i R^n \to \ldots \to R^n \to R
$$
with maps given by
$$
e_{j_1} \wedge \ldots \wedge e_{j_i}
\longmapsto
\sum_{a = 1}^i (-1)^{a + 1} x_{j_a} e_{j_1} \wedge \ldots
\wedge \hat e_{j_a} \wedge \ldots \wedge e_{j_i}
$$
It is easy to see that this is a complex $K_{\bullet}(R, x_{\bullet})$.
Note that the cokernel of the last map of $K_{\bullet}(R, x_{\bullet})$
is $\kappa$ by Lemma \ref{lemma-NAK} part (8).

\medskip\noindent
If $\kappa$ has finite projective dimension $d$, then we can find
a resolution $F_{\bullet} \to \kappa$ by finite free $R$-modules
of length $d$
(Lemma \ref{lemma-what-kind-of-resolutions-Noetherian-local}).
By Lemma \ref{lemma-add-trivial-complex}
we may assume all the maps in the complex $F_{\bullet}$
have the property that $\Im(F_i \to F_{i-1})
\subset \mathfrak m F_{i-1}$, because removing a trivial
summand from the resolution can at worst shorten the resolution.
By Lemma \ref{lemma-compare-resolutions} we can find a map
of complexes $\alpha : K_{\bullet}(R, x_{\bullet}) \to F_{\bullet}$
inducing the identity on $\kappa$. We will prove by induction
that the maps $\alpha_i : \wedge^i R^n = K_i(R, x_{\bullet}) \to F_i$
have the property that
$\alpha_i \otimes \kappa : \wedge^i \kappa^n \to F_i \otimes \kappa$
are injective. This shows that $F_n \not = 0$ and hence $d \geq n$
as desired.

\medskip\noindent
The result is clear for $i = 0$ because the composition
$R \xrightarrow{\alpha_0} F_0 \to \kappa$ is nonzero.
Note that $F_0$ must have rank $1$ since
otherwise the map $F_1 \to F_0$ whose cokernel is a single
copy of $\kappa$ cannot have image contained in $\mathfrak m F_0$.

\medskip\noindent
Next we check the case $i = 1$ as we feel that it is instructive;
the reader can skip this as the induction step will deduce the $i = 1$
case from the case $i = 0$. We saw above that
$F_0 = R$ and $F_1 \to F_0 = R$ has image $\mathfrak m$.
We have a commutative diagram
$$
\begin{matrix}
R^n & = & K_1(R, x_{\bullet}) & \to & K_0(R, x_{\bullet}) & = & R \\
& & \downarrow & & \downarrow & & \downarrow \\
& & F_1 & \to & F_0 & = & R
\end{matrix}
$$
where the rightmost vertical arrow is given by multiplication
by a unit. Hence we see that the image of the composition
$R^n \to F_1 \to F_0 = R$ is also equal to $\mathfrak m$.
Thus the map $R^n \otimes \kappa \to F_1 \otimes \kappa$
has to be injective since $\dim_\kappa (\mathfrak m / \mathfrak m^2) = n$.

\medskip\noindent
Let $i \geq 1$ and assume injectivity of $\alpha_j \otimes \kappa$ has been
proved for all $j \leq i - 1$. Consider the commutative diagram
$$
\begin{matrix}
\wedge^i R^n & = & K_i(R, x_{\bullet}) & \to & K_{i-1}(R, x_{\bullet})
& = & \wedge^{i-1} R^n \\
& & \downarrow & & \downarrow & & \\
& & F_i & \to & F_{i-1} & &
\end{matrix}
$$
We know that $\wedge^{i-1} \kappa^n \to F_{i-1} \otimes \kappa$
is injective. This proves that
$\wedge^{i-1} \kappa^n \otimes_{\kappa} \mathfrak m/\mathfrak m^2
\to F_{i-1} \otimes \mathfrak m/\mathfrak m^2$ is injective.
Also, by our choice of the complex, $F_i$ maps into
$\mathfrak mF_{i-1}$, and similarly for the Koszul complex.
Hence we get a commutative diagram
$$
\begin{matrix}
\wedge^i \kappa^n & \to &
\wedge^{i-1} \kappa^n \otimes \mathfrak m/\mathfrak m^2 \\
\downarrow & & \downarrow \\
F_i \otimes \kappa & \to & F_{i-1} \otimes \mathfrak m/\mathfrak m^2
\end{matrix}
$$
At this point it suffices to verify the map
$\wedge^i \kappa^n \to
\wedge^{i-1} \kappa^n \otimes \mathfrak m/\mathfrak m^2$
is injective, which can be done by hand.
\end{proof}

\begin{lemma}
\label{lemma-dim-gl-dim}
Let $R$ be a Noetherian local ring.
Suppose that the residue field $\kappa$ has finite
projective dimension $n$ over $R$.
In this case $\dim(R) \geq n$.
\end{lemma}

\begin{proof}
Let $F_{\bullet}$ be a finite resolution of $\kappa$ by finite free
$R$-modules (Lemma \ref{lemma-what-kind-of-resolutions-Noetherian-local}).
By Lemma \ref{lemma-add-trivial-complex}
we may assume all the maps in the complex $F_{\bullet}$
have to property that $\Im(F_i \to F_{i-1})
\subset \mathfrak m F_{i-1}$, because removing a trivial
summand from the resolution can at worst shorten the resolution.
Say $F_n \not = 0$ and $F_i = 0$ for $i > n$, so that
the projective dimension of $\kappa$ is $n$.
By Proposition \ref{proposition-what-exact} we see that
$\text{depth}_{I(\varphi_n)}(R) \geq n$ since $I(\varphi_n)$
cannot equal $R$ by our choice of the complex.
Thus by Lemma \ref{lemma-bound-depth} also $\dim(R) \geq n$.
\end{proof}

\begin{proposition}
\label{proposition-finite-gl-dim-regular}
Let $(R, \mathfrak m, \kappa)$ be a Noetherian local ring.
The following are equivalent
\begin{enumerate}
\item $\kappa$ has finite projective dimension as an $R$-module,
\item $R$ has finite global dimension,
\item $R$ is a regular local ring.
\end{enumerate}
Moreover, in this case the global dimension of $R$ equals
$\dim(R) = \dim_\kappa(\mathfrak m/\mathfrak m^2)$.
\end{proposition}

\begin{proof}
We have (3) $\Rightarrow$ (2) by
Proposition \ref{proposition-regular-finite-gl-dim}.
The implication (2) $\Rightarrow$ (1) is trivial.
Assume (1). By Lemmas \ref{lemma-length-resolution-residue-field}
and \ref{lemma-dim-gl-dim} we see that
$\dim(R) \geq \dim_\kappa(\mathfrak m /\mathfrak m^2)$.
Thus $R$ is regular, see Definition
\ref{definition-regular-local} and the discussion preceding it.
Assume the equivalent conditions (1) -- (3) hold.
By Proposition \ref{proposition-regular-finite-gl-dim}
the global dimension of $R$ is at most $\dim(R)$ and by
Lemma \ref{lemma-length-resolution-residue-field}
it is at least $\dim_\kappa(\mathfrak m/\mathfrak m^2)$.
Thus the stated equality holds.
\end{proof}

\begin{lemma}
\label{lemma-localization-of-regular-local-is-regular}
A Noetherian local ring $R$ is a regular local ring if and only if
it has finite global dimension. In this case
$R_{\mathfrak p}$ is a regular local ring for all primes $\mathfrak p$.
\end{lemma}

\begin{proof}
By Propositions \ref{proposition-finite-gl-dim-regular} and
\ref{proposition-regular-finite-gl-dim}
we see that a Noetherian local ring is a regular local ring if and only if
it has finite global dimension. Furthermore, any localization
$R_{\mathfrak p}$ has finite global dimension,
see Lemma \ref{lemma-localize-finite-gl-dim},
and hence is a regular local ring.
\end{proof}

\noindent
By Lemma \ref{lemma-localization-of-regular-local-is-regular}
it makes sense to make the following definition,
because it does not conflict with the earlier
definition of a regular local ring.

\begin{definition}
\label{definition-regular}
A Noetherian ring $R$ is said to be {\it regular}
if all the localizations $R_{\mathfrak p}$ at primes are
regular local rings.
\end{definition}

\noindent
It is enough to require the local rings at maximal ideals to be regular.
Note that this is not the same as asking $R$ to have finite
global dimension, even assuming $R$ is Noetherian. This is
because there is an example of a regular Noetherian ring
which does not have finite global dimension, namely because
it does not have finite dimension.

\begin{lemma}
\label{lemma-finite-gl-dim-finite-dim-regular}
Let $R$ be a Noetherian ring.
The following are equivalent:
\begin{enumerate}
\item $R$ has finite global dimension $n$,
\item $R$ is a regular ring of dimension $n$,
\item there exists an integer $n$ such that
all the localizations $R_{\mathfrak m}$ at maximal ideals
are regular of dimension $\leq n$ with equality for at least
one $\mathfrak m$, and
\item there exists an integer $n$ such that
all the localizations $R_{\mathfrak p}$ at prime ideals
are regular of dimension $\leq n$ with equality for at least
one $\mathfrak p$.
\end{enumerate}
\end{lemma}

\begin{proof}
This follows from the discussion above. More precisely, it follows
by combining Definition \ref{definition-regular} with
Lemma \ref{lemma-finite-gl-dim-primes}
and Proposition \ref{proposition-finite-gl-dim-regular}.
\end{proof}

\begin{lemma}
\label{lemma-flat-under-regular}
Let $R \to S$ be a local homomorphism of local Noetherian rings.
Assume that $R \to S$ is flat and that $S$ is regular.
Then $R$ is regular.
\end{lemma}

\begin{proof}
Let $\mathfrak m \subset R$ be the maximal ideal
and let $\kappa = R/\mathfrak m$ be the residue field.
Let $d = \dim S$.
Choose any resolution $F_\bullet \to \kappa$
with each $F_i$ a finite free $R$-module. Set
$K_d = \Ker(F_{d - 1} \to F_{d - 2})$.
By flatness of $R \to S$ the complex
$0 \to K_d \otimes_R S \to F_{d - 1} \otimes_R S \to \ldots
\to F_0 \otimes_R S \to \kappa \otimes_R S \to 0$
is still exact. Because the global dimension of $S$
is $d$, see Proposition \ref{proposition-regular-finite-gl-dim},
we see that $K_d \otimes_R S$ is a finite free $S$-module
(see also Lemma \ref{lemma-independent-resolution}).
By Lemma \ref{lemma-finite-projective-descends} we see
that $K_d$ is a finite free $R$-module.
Hence $\kappa$ has finite projective dimension and $R$ is regular by
Proposition \ref{proposition-finite-gl-dim-regular}.
\end{proof}








